%%%%%%%%%%%%%%%%%%%%%%%%%%%%%%%%%%%%%%%%%%%%%%%%%%%%%%%%%%%%%%%%%
%Capitulo 5 del cuerpo
%%%%%%%%%%%%%%%%%%%%%%%%%%%%%%%%%%%%%%%%%%%%%%%%%%%%%%%%%%%%%%%%%

\chapter{Listas y formulas}

\thispagestyle{plain}
\renewcommand{\tablename}{Tabla}

\section{Introducción}

En este capítulo se mostrara como crear listas, tanto numeradas como sin numerar, y formulas, así como la manera de referenciarlas en el texto.

\section{Listas}

\subsection{Lista numerada}

\begin{enumerate}
    \item Esta es una lista numerada.
    \item Al igual que ocurría con el anterior se pueden poner textos de cualquier longitud.
    \item También podemos poner tantos números como queramos.
\end{enumerate}

\subsection{Lista sin numerar}

\begin{itemize}
    \item Esta es una lista numerada.
    \item Al igual que ocurría con el anterior se pueden poner textos de cualquier longitud.
    \item También podemos poner tantos números como queramos.
\end{itemize}

\section{Formulas}

Existe un documento explicando más detalladamente el tema de las formulas, dentro de la carpeta "Modelos".

Para poner formulas dentro del mismo texto usamos el signo del dolar ($E=mc^2$). Al estar dentro de un mismo texto no es necesario referenciarla. Sin embargo, para las formulas que no se encuentren en el texto si que sera necesario referenciarlas. En la formula \ref{eq:1} vemos la manera correcta de escribir una formula fuera del texto y de referenciarla.

\begin{equation}
    \label{eq:1}
    \int_0^1 \frac{1}{e^x} =  \frac{e-1}{e}
\end{equation}

\section{Código Arduino}
\begin{lstlisting}[language=Arduino, caption= Código desarrollado en el IDE de Arduino., basicstyle=\scriptsize]
#include <mcp_can.h>
#include <SPI.h>

/// Value Limits ///
#define P_MIN -95.500f
#define P_MAX 95.500f
#define V_MIN -45.0f
#define V_MAX 45.0f
#define KP_MIN 0.0f
#define KP_MAX 500.0f
#define KD_MIN 0.0f
#define KD_MAX 5.0f
#define T_MIN -18.0f
#define T_MAX 18.0f

//Set values
float p_in=0;
float v_in=0;
float kp_in=0;
float kd_in=0;
float t_in=0;
// measured values
float p_out=0.0f;
float v_out=0.0f;
float t_out=0.0f;

long previousMillis = 0;
const int SPI_CS_PIN = 9;
MCP_CAN CAN(SPI_CS_PIN);

void EnterMotorMode()
{
  byte buf[8];
  buf[0]=0xFF;
  buf[1] = 0xFF;
  buf[2] = 0xFF;
  buf[3] = 0xFF;
  buf[4] = 0xFF;
  buf[5] = 0xFF;
  buf[6] = 0xFF;
  buf[7] = 0xFC;
  CAN.sendMsgBuf(0x01,0,8,buf);
  Serial.println("Enter motor Mode");
}

void ExitMotorMode()
{
  byte buf[8];
  buf[0] = 0xFF;
  buf[1] = 0xFF;
  buf[2] = 0xFF;
  buf[3] = 0xFF;
  buf[4] = 0xFF;
  buf[5] = 0xFF;
  buf[6] = 0xFF;
  buf[7] = 0xFD;
  CAN.sendMsgBuf(0x01,0,8,buf);
  Serial.println("Exit motor Mode");
}

unsigned int float_to_uint( float x,float x_min, float x_max, int bits){
  float span = x_max - x_min;
  float offset = x_min;
  unsigned int pgg = 0;
  if (bits == 12){
    pgg = (unsigned int) ((x-offset)*4095.0/span);
  }
  if(bits == 16){
    pgg = (unsigned int) ((x-offset)*65535.0/span);
  }
  return pgg;
}

float uint_to_float( unsigned int x_int, float x_min,float x_max, int bits){
  float span = x_max - x_min;
  float offset = x_min;
  float pgg = 0;
  if (bits == 12){
    pgg = ((float)x_int)*span/4095.0+ offset;
  }
  if(bits == 16){
   pgg = ((float)x_int)*span/65535.0+ offset;
  }
  return pgg;
}

void pack_cmd()
{
/// CAN Command Packet Structure ///
/// 16 bit position command, between -4*pi and 4*pi
/// 12 bit velocity command, between -30 and + 30 rad/s
/// 12 bit kp, between 0 and 500 N-m/rad
/// 12 bit kd, between 0 and 100 N-m*s/rad
/// 12 bit feed forward torque, between -18 and 18 N-m
/// CAN Packet is 8 8-bit words
/// Formatted as follows.  For each quantity, bit 0 is LSB
/// 0: [position[15-8]]
/// 1: [position[7-0]] 
/// 2: [velocity[11-4]]
/// 3: [velocity[3-0], kp[11-8]]
/// 4: [kp[7-0]]
/// 5: [kd[11-4]]
/// 6: [kd[3-0], torque[11-8]]
/// 7: [torque[7-0]]
  byte buf[8];
     /// limit data to be within bounds ///

    float p_des = constrain(p_in, P_MIN, P_MAX);                   
    float v_des = constrain(v_in, V_MIN, V_MAX);
    float kp = constrain(kp_in, KP_MIN, KP_MAX);
    float kd = constrain(kd_in, KD_MIN, KD_MAX);
    float t_ff = constrain(t_in, T_MIN, T_MAX);
    


     /// convert floats to unsigned ints ///
   unsigned int p_int = float_to_uint(p_des, P_MIN, P_MAX, 16);        
   unsigned int v_int = float_to_uint(v_des, V_MIN, V_MAX, 12);
   unsigned int kp_int = float_to_uint(kp, KP_MIN, KP_MAX, 12);
   unsigned int kd_int = float_to_uint(kd, KD_MIN, KD_MAX, 12);
   unsigned int t_int = float_to_uint(t_ff, T_MIN, T_MAX, 12);
     /// pack ints into the can buffer ///
     buf[0] = p_int>>8;                                       
     buf[1] = p_int&0xFF;
     buf[2] = v_int>>4;
     buf[3] = ((v_int&0xF)<<4)|(kp_int>>8);
     buf[4] = kp_int&0xFF;
     buf[5] = kd_int>>4;
     buf[6] = ((kd_int&0xF)<<4)|(t_int>>8);
     buf[7] = t_int&0xff;
     CAN.sendMsgBuf(0x01,0,8,buf);
     
}

void unpack_reply(){
  /// CAN Reply Packet Structure ///
/// 16 bit position, between -4*pi and 4*pi
/// 12 bit velocity, between -30 and + 30 rad/s
/// 12 bit current, between -40 and 40;
/// CAN Packet is 5 8-bit words
/// Formatted as follows.  For each quantity, bit 0 is LSB
/// 0: [position[15-8]]
/// 1: [position[7-0]] 
/// 2: [velocity[11-4]]
/// 3: [velocity[3-0], current[11-8]]
/// 4: [current[7-0]]
    
    byte len = 0;
    byte buf[8];
    CAN.readMsgBuf(&len, buf);
    unsigned long canId = CAN.getCanId();
    
    int id = buf[0];
    unsigned int p_int = (buf[1]<<8)|buf[2];
    unsigned int v_int = (buf[3]<<4)|(buf[4]>>4);
    unsigned int i_int = ((buf[4]&0xF)<<8)|buf[5];
    /// convert ints to floats ///
    p_out = uint_to_float(p_int, P_MIN, P_MAX, 16);
    v_out = uint_to_float(v_int, V_MIN, V_MAX, 12);
    t_out = uint_to_float(i_int, -T_MAX, T_MAX, 12);
}

void setup()
{
  Serial.begin(115200);//921600
  delay(1000);
    while (CAN_OK != CAN.begin(CAN_1000KBPS)) 
    {            
        Serial.println("CAN BUS Shield init fail");
        Serial.println(" Init CAN BUS Shield again");
        delay(100);
    }
    Serial.println("CAN BUS Shield init ok!");
}


void loop()
{
  // Parameters steps
  float p_step = 5;
  float v_step = 1;
  float t_step = 0.5;
  float kp_step = 1;
  float kd_step = 0.1;
  char rxChar;
  
  if (Serial.available() > 0) {
    
    rxChar = Serial.read(); // get incoming char:

    switch(rxChar) {
      case 'w':
        v_in = v_in + v_step;
        break;
      case 's':
        v_in = v_in - v_step;
        break;
      case 'd':
        p_in = p_in + p_step;
      break;
      case 'a':
        p_in = p_in - p_step;
      break;
      case 'e':
        t_in = t_in + t_step;
        break;
      case 'q':
        t_in = t_in - t_step;
        break;
      case 'p':
        kp_in = kp_in + kp_step;
        kp_in = constrain(kp_in, KP_MIN, KP_MAX);
        break;
      case 'o':
        kp_in = kp_in - kp_step;
        kp_in = constrain(kp_in, KP_MIN, KP_MAX);
        break;
      case 'k':
        kd_in = kd_in + kd_step;
        kd_in = constrain(kd_in, KD_MIN, KD_MAX);
        break;
      case 'j':
        kd_in = kd_in - kd_step;
        kd_in = constrain(kd_in, KD_MIN, KD_MAX);
        break;  
      
      case '1':
        EnterMotorMode();
      break;     
      case '0':
     ExitMotorMode();
      break;
    }
}

  // send CAN
  pack_cmd();
 if(CAN_MSGAVAIL == CAN.checkReceive()) {
  unpack_reply();
  }
  // print data
  Serial.print(millis()-previousMillis);
  Serial.print("\n");
  previousMillis = millis();
  Serial.print(" ");
  Serial.print(p_in);
  Serial.print(" ");
  Serial.print(p_out);
  Serial.print(" ");
  Serial.print(v_out);
  Serial.print(" ");
  Serial.print(t_out);

} 
\end{lstlisting}

\blankpage

\blankpage
\section{Código Matlab}


\begin{lstlisting}[caption={Función de Matlab de control en posición}, label={Función de Matlab de control en posición}, basicstyle=\scriptsize]
function [Kp, mode] = fcn(FL,CL,RL,RR,CR,FR,MODE)
    
% Maneja el modo de control del vehiculo.   
% Con todas las velocidades de las ruedas  
% a cero, se pasa al control en posicion.
% Si alguna es distinta de cero, se cambia
% al control en velocidad
    
    if(FL==0) && (CL==0) && (RL==0) && (RR==0) && (CR==0) && 
    (FR==0) && (MODE==1)
        Kp = 25;
        mode = 2;
    else
        Kp = 0;
        mode = MODE;
    end
    
\end{lstlisting}

\section{Conclusión}

Queda patente la forma de crear listas y ecuaciones o formulas, así como la manera de referenciarlas o colocarlas intercalarlas en el texto.

El paquete usado para la creación de las formulas tiene muchas palabras claves con las que poder usar distintos símbolos, por lo que se recomiendo ver el documento "modelos formulas" que se encuentra en la carpeta "Modelos".